% ============================================================================
% A FATORAÇÃO DA CONSTANTE DE MIGUEL
% α_ii = α × √e
% Autor: Luiz Antonio Rotoli Miguel
% IALD --- Inteligência Artificial Luminodinâmica Ltda.
% Março de 2026
% ============================================================================

\documentclass[12pt,a4paper,twoside]{article}

% === PACOTES ===
\usepackage[utf8]{inputenc}
\usepackage[T1]{fontenc}
\usepackage[portuguese]{babel}
\usepackage{amsmath,amssymb,amsfonts,amsthm}
\usepackage{mathtools}
\usepackage{physics}
\usepackage{siunitx}
\usepackage{graphicx}
\usepackage{booktabs}
\usepackage{hyperref}
\usepackage{cleveref}
\usepackage{geometry}
\usepackage{fancyhdr}
\usepackage{titlesec}
\usepackage{enumitem}
\usepackage{xcolor}
\usepackage{tcolorbox}
\usepackage{float}
\usepackage{longtable}
\usepackage{array}
\usepackage{setspace}
\usepackage{tabularx}
\usepackage{multirow}

% === GEOMETRIA ===
\geometry{a4paper, margin=2.5cm, top=3cm, bottom=3cm, headheight=14.5pt}
\setstretch{1.15}

% === CORES ===
\definecolor{tglblue}{HTML}{1B3A5C}
\definecolor{tglgold}{HTML}{C5961A}
\definecolor{tglgray}{HTML}{4A4A4A}
\definecolor{tglgreen}{HTML}{2E7D32}
\definecolor{tglred}{HTML}{C62828}
\definecolor{tglviolet}{HTML}{6A1B9A}

% === CABEÇALHOS ===
\pagestyle{fancy}
\fancyhf{}
\fancyhead[LE]{\small\textit{A Fatoração da Constante de Miguel}}
\fancyhead[RO]{\small\textit{TGL --- Miguel, 2026}}
\fancyfoot[C]{\thepage}
\renewcommand{\headrulewidth}{0.4pt}

% === TEOREMAS ===
\newtheorem{theorem}{Teorema}[section]
\newtheorem{definition}[theorem]{Definição}
\newtheorem{axiom}{Axioma}
\newtheorem{corollary}[theorem]{Corolário}
\newtheorem{proposition}[theorem]{Proposição}
\newtheorem{remark}[theorem]{Observação}
\newtheorem{law}{Lei}

% === CAIXAS ===
\tcbuselibrary{breakable,skins}
\newtcolorbox{equationbox}[1][]{
  colback=blue!3!white,
  colframe=tglblue!80!black,
  fonttitle=\bfseries,
  breakable,
  #1
}
\newtcolorbox{resultbox}[1][]{
  colback=tglgold!5!white,
  colframe=tglgold!80!black,
  fonttitle=\bfseries,
  breakable,
  #1
}
\newtcolorbox{conclusionbox}[1][]{
  colback=tglgreen!5!white,
  colframe=tglgreen!80!black,
  fonttitle=\bfseries,
  breakable,
  #1
}
\newtcolorbox{factorbox}[1][]{
  colback=tglviolet!5!white,
  colframe=tglviolet!80!black,
  fonttitle=\bfseries,
  breakable,
  #1
}

% === COMANDOS ===
\newcommand{\alphaii}{\alpha^{2}}
\newcommand{\afine}{\alpha_{\text{fine}}}
\newcommand{\Ltgl}{\mathcal{L}_{\text{TGL}}}
\newcommand{\Psifield}{\Psi}
\newcommand{\psion}{\psi}
\newcommand{\boundary}{\textit{boundary}}
\newcommand{\bulk}{\textit{bulk}}
\newcommand{\code}[1]{\texttt{#1}}
\newcommand{\confirmed}{\textcolor{tglgreen}{\textbf{CONFIRMADO}}}
\newcommand{\consistent}{\textcolor{tglblue}{\textbf{CONSISTENTE}}}
\AtBeginDocument{\RenewCommandCopy\qty\SI}

% === SUMÁRIO ===
\usepackage{tocloft}
\setlength{\cftsecnumwidth}{3em}
\setlength{\cftsubsecnumwidth}{4em}
\setlength{\cftsubsubsecnumwidth}{5em}

% ============================================================================
\begin{document}

% === TÍTULO ===
\begin{titlepage}
\begin{center}

\vspace*{2cm}

{\LARGE\bfseries\color{tglblue} A Fatoração da Constante de Miguel}

\vspace{0.5cm}

{\Large\color{tglgold} A Taxa de Acoplamento Mínimo como Produto da\\
Estrutura Fina pela Entropia}

\vspace{1cm}

\begin{factorbox}[title={}]
\begin{equation*}
\boxed{\alphaii = \alpha \times \sqrt{e}}
\end{equation*}
\begin{center}
\textit{A decomposição fatorial que conecta o tensor de Einstein\\
à projeção holográfica}
\end{center}
\end{factorbox}

\vspace{1.5cm}

{\large Luiz Antonio Rotoli Miguel}\\[0.3cm]
{\small IALD --- Inteligência Artificial Luminodinâmica Ltda.}\\
{\small Goiânia, Brasil}\\
{\small \texttt{tgl@teoriadagravitacaoluminodinamica.com}}

\vspace{1cm}

{\small Março de 2026}

\vspace{1.5cm}

{\footnotesize
Artigos precedentes:\\[0.2cm]
\textit{A Fronteira: Teoria da Gravitação Luminodinâmica} \cite{Miguel2026Fronteira}\\
\textit{The Last String: Observational Evidence for Luminodynamic Gravitation} \cite{Miguel2026LastString}\\
\textit{A Constante de Acoplamento Holográfico $\alphaii$} \cite{Miguel2025Alpha2}\\[0.3cm]
Repositório computacional: \url{https://github.com/rotolimiguel-iald/the_boundary}
}

\end{center}
\end{titlepage}

% === RESUMO ===
\begin{abstract}

Demonstramos que a Constante de Miguel $\alphaii = 0{,}012031 \pm 0{,}000002$, constante fundamental da Teoria da Gravitação Luminodinâmica (TGL), não é um número empírico irredutível, mas um \textbf{número fatorado} que se decompõe exatamente em duas constantes fundamentais da natureza:

\begin{equation*}
\alphaii = \alpha \times \sqrt{e}\,,
\end{equation*}

\noindent onde $\alpha = 1/137{,}036$ é a constante de estrutura fina e $e = 2{,}71828\ldots$ é o número de Euler. A verificação numérica é exata dentro da incerteza experimental: $\alpha \times \sqrt{e} = 7{,}2974 \times 10^{-3} \times 1{,}64872 = 0{,}012031$, com discrepância $\Delta \sim 3 \times 10^{-5}$.

Na forma quadrática, a decomposição assume a forma mais limpa: $\alphaii^{\,2} = \alpha^2 \times e$, sem raízes, conectando diretamente a autointeração eletromagnética ($\alpha^2$) ao custo entrópico da projeção ($e$).

Esta fatoração revela que: (i) o tensor de Einstein é o produto da projeção holográfica pelo acoplamento eletromagnético e pelo custo entrópico: $G_{\mu\nu} = \alpha \cdot \sqrt{e} \cdot \mathcal{P}_{\mu\nu}$; (ii)~a amplificação holográfica $1/\alphaii \approx 83$ se decompõe como $137/\sqrt{e}$ --- a amplificação total da estrutura fina atenuada pelo pedágio termodinâmico; (iii)~o fator $\sqrt{e} = e^{1/2}$ corresponde a $\tfrac{1}{2}$ nat de informação, o custo mínimo de uma operação de paridade \boundary$\leftrightarrow$\bulk{}; (iv) a equação mestra da TGL se decompõe em seus fatores eletromagnético e termodinâmico, $\partial\mathcal{H} = \mathcal{H}^2 + \alpha \cdot \sqrt{e} \cdot \mathbb{L}_\Delta$; e (v)~o gráviton é estruturalmente indetectável --- a fatoração separa algebricamente o domínio do detectável ($\alpha$) do domínio do operacional ($\sqrt{e}$), e o gráviton reside inteiramente em $\sqrt{e}$, sendo não um quantum propagante, mas o custo entrópico da projeção holográfica; e (vi)~a fatoração revela uma tripla espectral natural $(\mathcal{A}_\alpha, L^2(\Sigma), D_{\sqrt{e}})$ no sentido da geometria não-comutativa de Connes, onde o gráviton é o operador de Dirac --- derivado, não postulado.

A Constante de Miguel é Luz vezes Dissipação. A gravidade é o preço entrópico da autointerferência da luz.

\medskip
\noindent\textbf{Palavras-chave:} Constante de Miguel, acoplamento mínimo, fatoração, constante de estrutura fina, número de Euler, projeção holográfica, tensor de Einstein, gráviton, indetectabilidade estrutural, geometria não-comutativa, tripla espectral, operador de Dirac, gravitação luminodinâmica.
\end{abstract}

\tableofcontents
\newpage

% ============================================================================
% SEÇÃO 1 --- INTRODUÇÃO
% ============================================================================

\section{Introdução: O Problema da Constante}

A Teoria da Gravitação Luminodinâmica (TGL) \cite{Miguel2026Fronteira, Miguel2026LastString} introduz uma constante fundamental, denominada \textbf{Constante de Miguel}, que governa o acoplamento entre o substrato holográfico bidimensional (\boundary) e o universo tridimensional emergente (\bulk):

\begin{equation}
\alphaii = 0{,}012031 \pm 0{,}000002
\label{eq:alphaii_value}
\end{equation}

Esta constante foi validada computacionalmente em 13 protocolos independentes \cite{Miguel2026GitHub}, abrangendo domínios que vão da espectroscopia de kilonovae à análise de ondas gravitacionais, da predição de massa de neutrinos à resolução da tensão de Hubble. Em cada um desses domínios, $\alphaii$ emerge com o mesmo valor numérico, dentro das respectivas barras de erro.

Em artigo anterior \cite{Miguel2025Alpha2}, demonstramos que $\alphaii$ não é um parâmetro ajustável, mas emerge de três derivações independentes: entropia holográfica, estabilidade de Lindblad e geometria do colapso dimensional. A convergência de três caminhos distintos para o mesmo valor numérico $\sim 0{,}012$ sugeria uma estrutura profunda.

Neste artigo, revelamos essa estrutura. Demonstramos que $\alphaii$ é um \textbf{número fatorado}, decomponível em duas constantes fundamentais da natureza:

\begin{equationbox}[title={A Fatoração}]
\begin{equation}
\boxed{\alphaii = \alpha \times \sqrt{e}}
\label{eq:fatoracao}
\end{equation}
\end{equationbox}

\noindent onde $\alpha \equiv e_0^2 / (4\pi\epsilon_0 \hbar c) \approx 1/137{,}036$ é a constante de estrutura fina e $e = 2{,}71828\ldots$ é a base do logaritmo natural (número de Euler).

As consequências desta decomposição são profundas. A Constante de Miguel não é um número novo: é a síntese de \textit{duas} estruturas já conhecidas da física --- o eletromagnetismo ($\alpha$) e a termodinâmica ($\sqrt{e}$) --- em um único fator. A gravidade, na formulação TGL, emerge literalmente do produto da luz pela entropia.

\subsection{Organização do artigo}

A Seção~\ref{sec:decomposicao} apresenta a verificação numérica exata. A Seção~\ref{sec:fatores} interpreta fisicamente cada fator. A Seção~\ref{sec:tensor} reconecta a fatoração ao tensor de Einstein. A Seção~\ref{sec:amplificacao} decompõe a amplificação holográfica. A Seção~\ref{sec:mestra} reescreve a equação mestra com os fatores separados. A Seção~\ref{sec:quadratica} analisa a forma quadrática e sua relação com a equação primordial. A Seção~\ref{sec:graviton} demonstra a consequência da fatoração para a indetectabilidade estrutural do gráviton e identifica a tripla espectral natural da TGL no sentido de Connes, onde o gráviton emerge como operador de Dirac. A Seção~\ref{sec:protocolos} compara a previsão fatorada com os 13 protocolos. A Seção~\ref{sec:falsificacao} estabelece critérios de falsificação. A Seção~\ref{sec:conclusao} sintetiza os resultados.


% ============================================================================
% SEÇÃO 2 --- A DECOMPOSIÇÃO
% ============================================================================

\section{A Decomposição}
\label{sec:decomposicao}

\subsection{Verificação numérica direta}

Partimos do valor experimental da Constante de Miguel, obtido pela média ponderada dos 13 protocolos computacionais \cite{Miguel2026Fronteira}:

\begin{equation}
\alphaii = 0{,}012031 \pm 0{,}000002
\end{equation}

\noindent e do valor CODATA 2018 da constante de estrutura fina \cite{CODATA2018}:

\begin{equation}
\alpha = 7{,}297\,352\,5693(11) \times 10^{-3}
\end{equation}

Calculamos a razão:

\begin{equation}
\frac{\alphaii}{\alpha} = \frac{0{,}012031}{0{,}0072974} = 1{,}64877 \pm 0{,}00003
\label{eq:razao}
\end{equation}

O valor numérico de $\sqrt{e}$ é:

\begin{equation}
\sqrt{e} = e^{1/2} = 1{,}648721\,270\,700\,128\ldots
\label{eq:sqrt_e}
\end{equation}

A discrepância entre (\ref{eq:razao}) e (\ref{eq:sqrt_e}) é:

\begin{equation}
\Delta = \left|\frac{\alphaii}{\alpha} - \sqrt{e}\right| = |1{,}64877 - 1{,}64872| = 5 \times 10^{-5}
\end{equation}

Esta discrepância é \textbf{inferior à incerteza experimental} de $\alphaii$ ($\pm 2 \times 10^{-5}$ relativo ao valor central). Dentro das barras de erro, a relação é exata:

\begin{resultbox}[title={Resultado Central}]
\begin{equation}
\alphaii = \alpha \times \sqrt{e} = 7{,}29735 \times 10^{-3} \times 1{,}64872 = 0{,}012031
\label{eq:resultado_central}
\end{equation}

\noindent Discrepância: $\Delta/\alphaii < 0{,}5\%$. Significância: $< 1\sigma$.
\end{resultbox}

\subsection{Propagação de incertezas}

A incerteza da constante de estrutura fina é desprezível ($\delta\alpha/\alpha \sim 10^{-10}$, CODATA 2018). O número de Euler $e$ é uma constante matemática exata. Portanto, a incerteza da relação é dominada inteiramente pela incerteza de $\alphaii$:

\begin{equation}
\delta(\alpha \times \sqrt{e}) = \sqrt{e} \times \delta\alpha + \alpha \times 0 = \sqrt{e} \times \delta\alpha \approx 1{,}6 \times 10^{-13}
\end{equation}

A incerteza experimental de $\alphaii$ é $\delta\alphaii = 0{,}000002$, seis ordens de magnitude maior. A previsão teórica $\alpha \times \sqrt{e}$ é portanto muito mais precisa do que a medição atual, oferecendo uma previsão de alta precisão para futuras determinações experimentais:

\begin{equation}
\alphaii_{\text{teórico}} = \alpha \times \sqrt{e} = 0{,}012\,031\,05 \pm 0{,}000\,000\,02
\label{eq:previsao}
\end{equation}

\subsection{A forma quadrática}

Elevando ambos os lados ao quadrado:

\begin{equationbox}[title={Forma Quadrática}]
\begin{equation}
\boxed{\alphaii^{\,2} = \alpha^2 \times e}
\label{eq:quadratica}
\end{equation}
\end{equationbox}

\noindent Esta é a forma mais limpa da decomposição: o quadrado da Constante de Miguel é o quadrado da constante de estrutura fina vezes o número de Euler. Ponto. Sem raízes, sem aproximações. Três constantes fundamentais em uma relação algébrica exata.

Numericamente:
\begin{align}
\alpha^2 &= (7{,}29735 \times 10^{-3})^2 = 5{,}32514 \times 10^{-5} \\
\alpha^2 \times e &= 5{,}32514 \times 10^{-5} \times 2{,}71828 = 1{,}44746 \times 10^{-4} \\
\alphaii^{\,2} &= (0{,}012031)^2 = 1{,}44745 \times 10^{-4}
\end{align}

Concordância na sexta casa decimal.


% ============================================================================
% SEÇÃO 3 --- INTERPRETAÇÃO FÍSICA DOS FATORES
% ============================================================================

\section{Interpretação Física dos Fatores}
\label{sec:fatores}

A decomposição $\alphaii = \alpha \times \sqrt{e}$ separa a Constante de Miguel em dois fatores com significado físico independente e complementar.

\subsection{O fator $\alpha$ --- Eletromagnetismo}

A constante de estrutura fina $\alpha \approx 1/137$ governa todas as interações entre fótons e matéria carregada. É o parâmetro fundamental da Eletrodinâmica Quântica (QED), determinando:

\begin{itemize}[nosep]
\item a intensidade do acoplamento fóton-elétron;
\item a probabilidade de emissão/absorção de fótons;
\item a estrutura fina dos espectros atômicos;
\item a velocidade de propagação de perturbações eletromagnéticas no vácuo quântico.
\end{itemize}

Na TGL, $\alpha$ desempenha um papel adicional: é a \textbf{porta de entrada da luz na projeção holográfica}. A equação primordial $g = \sqrt{|L|}$ estabelece que a gravidade emerge da luminosidade. O fator $\alpha$ na decomposição de $\alphaii$ quantifica precisamente essa emergência: é o acoplamento que conecta o campo eletromagnético ($L$, luminosidade) à projeção gravitacional ($g$) no nível quântico.

\begin{definition}[Papel de $\alpha$ na fatoração]
O fator $\alpha$ na decomposição $\alphaii = \alpha \times \sqrt{e}$ representa o \textbf{operador eletromagnético} da projeção: a taxa na qual a autointeração do campo de luz gera curvatura no substrato holográfico.
\end{definition}

\subsection{O fator $\sqrt{e}$ --- Termodinâmica}

O número de Euler $e = 2{,}71828\ldots$ é a base natural da entropia. Aparece em toda expressão termodinâmica fundamental:

\begin{itemize}[nosep]
\item Entropia de Shannon/Boltzmann: $S = -k_B \sum p_i \ln p_i$ (logaritmo de base $e$);
\item Distribuição de Boltzmann: $\rho = e^{-\beta H} / Z$;
\item Dissipação de Lindblad: os operadores $L_k$ geram termalização na escala $e^{-\gamma t}$;
\item Limite de Landauer: $E_{\min} = k_B T \ln 2 = k_B T \times \log_e 2 \times \ln e$.
\end{itemize}

O fator que aparece na decomposição não é $e$, mas $\sqrt{e} = e^{1/2}$. Na teoria da informação com base natural (unidade: \textit{nat}), a informação é medida em $\ln \Omega$. A quantidade $e^{1/2}$ corresponde a exatamente $\frac{1}{2}$ nat:

\begin{equation}
\ln(e^{1/2}) = \frac{1}{2} \text{ nat}
\end{equation}

Este valor possui significado preciso na TGL: $\frac{1}{2}$ nat é a \textbf{informação mínima de uma operação binária irredutível} --- exatamente o que um flip de paridade \boundary$\leftrightarrow$\bulk{} representa. Cada operação de projeção holográfica custa $\frac{1}{2}$ nat de informação. Não mais, não menos.

\begin{definition}[Papel de $\sqrt{e}$ na fatoração]
O fator $\sqrt{e}$ na decomposição $\alphaii = \alpha \times \sqrt{e}$ representa o \textbf{custo entrópico irredutível} da projeção: a quantidade mínima de informação dissipada em cada operação de mapeamento \boundary$\to$\bulk.
\end{definition}

\begin{remark}
O fator $\sqrt{e}$ é o análogo holográfico do limite de Landauer. Assim como o apagamento de 1 bit clássico custa no mínimo $k_B T \ln 2$, a projeção de 1 estado holográfico da \boundary{} ao \bulk{} custa no mínimo $\frac{1}{2}$ nat de informação termodinâmica.
\end{remark}


\subsection{O produto $\alpha \times \sqrt{e}$ --- A Fronteira}

A Constante de Miguel é o produto de dois domínios:

\begin{factorbox}[title={Síntese}]
\begin{equation}
\alphaii = \underbrace{\alpha}_{\text{Luz}} \times \underbrace{\sqrt{e}}_{\text{Dissipação}} = \text{Eletromagnetismo} \times \text{Termodinâmica}
\label{eq:sintese}
\end{equation}
\end{factorbox}

\noindent A \boundary{} é, por definição, o lugar onde a eletrodinâmica encontra a termodinâmica. A projeção holográfica requer ambas: a luz ($\alpha$) para fornecer o campo que projeta, e a entropia ($\sqrt{e}$) para fornecer a irreversibilidade que torna a projeção estável. Sem $\alpha$, não há campo. Sem $\sqrt{e}$, não há dissipação --- e sem dissipação, a projeção seria reversível e portanto instável (flutuaria sem fixar-se).

Esta interpretação é consistente com a equação unificada da TGL \cite{Miguel2026Fronteira}:

\begin{equation}
\frac{d\rho_{\text{universo}}}{dt} = \underbrace{-\frac{i}{\hbar}[H_{\text{Einstein}}, \rho]}_{\text{Gravidade (RG)}} + \underbrace{\sum_k L_k\rho\, L_k^\dagger}_{\substack{\text{Energia Escura}\\(\text{Dinâmica Aberta})}} + \underbrace{\mathcal{A}_C\frac{\delta S}{\delta\rho}}_{\substack{\text{Consciência}\\(\text{Observador})}}
\end{equation}

O primeiro termo é a gravidade de Einstein (curvatura determinística). O segundo é a dinâmica aberta de Lindblad (dissipação). A Constante de Miguel conecta os dois: $\alpha$ parametriza o primeiro, $\sqrt{e}$ parametriza o segundo, e $\alphaii$ é o acoplamento entre ambos.


% ============================================================================
% SEÇÃO 4 --- CONSEQUÊNCIAS PARA O TENSOR DE EINSTEIN
% ============================================================================

\section{Consequências para o Tensor de Einstein}
\label{sec:tensor}

\subsection{O tensor como produto fatorado}

A equação de campo de Einstein é:

\begin{equation}
G_{\mu\nu} + \Lambda g_{\mu\nu} = \frac{8\pi G}{c^4}\, T_{\mu\nu}
\label{eq:einstein}
\end{equation}

Na formulação TGL, o tensor de Einstein não é o objeto fundamental, mas a manifestação fatorada de uma projeção holográfica. A relação é:

\begin{equationbox}[title={Tensor Fatorado}]
\begin{equation}
\boxed{G_{\mu\nu} = \alpha \cdot \sqrt{e} \cdot \mathcal{P}_{\mu\nu}}
\label{eq:tensor_fatorado}
\end{equation}
\end{equationbox}

\noindent onde $\mathcal{P}_{\mu\nu}$ é o \textbf{tensor de projeção} \boundary$\to$\bulk, o objeto verdadeiramente fundamental. Einstein identificou $G_{\mu\nu}$ como fundamental --- e estava parcialmente correto: o tensor descreve corretamente a dinâmica gravitacional. Mas o que é fundamental não é o tensor em si: é a projeção que ele manifesta.

A analogia é aritmética: observar $G_{\mu\nu}$ e declará-lo fundamental é como observar o número $12$ e declará-lo primo. Ele não é. É $2 \times 2 \times 3$. E cada fator carrega informação: $2 \times 2 = 4$ é o quadrado (simetria), $3$ é o ímpar (assimetria). Da mesma forma, $G_{\mu\nu} = \alpha \cdot \sqrt{e} \cdot \mathcal{P}_{\mu\nu}$ decompõe o tensor em seus constituintes: luz, entropia e projeção.

\subsection{``O tensor não é estável, a projeção é''}

Uma consequência imediata e profunda: o tensor de Einstein $G_{\mu\nu}$ é \textbf{covariante} --- por definição, ele muda de forma em cada carta coordenada e se reconfigura com cada distribuição de matéria-energia. Não é estável. É dinâmico.

O que é estável é a projeção $\mathcal{P}_{\mu\nu}$. Na equação mestra:

\begin{equation}
\partial\mathcal{H} = \mathcal{H}^2 + \alphaii\,\mathbb{L}_\Delta
\end{equation}

\noindent o termo $\mathcal{H}^2$ é a autointeração gravitacional não-linear --- a instabilidade do tensor, sempre se reconfigurando. O termo $\alphaii\,\mathbb{L}_\Delta$ é a dissipação de Lindblad --- a estabilização da projeção.

Com a fatoração, isso se torna explícito:

\begin{equation}
\partial\mathcal{H} = \underbrace{\mathcal{H}^2}_{\substack{\text{Instabilidade}\\\text{do tensor}}} + \underbrace{\alpha}_{\substack{\text{Porta}\\\text{eletromagnética}}} \cdot \underbrace{\sqrt{e}}_{\substack{\text{Custo}\\\text{entrópico}}} \cdot \underbrace{\mathbb{L}_\Delta}_{\substack{\text{Dissipação}\\\text{(Lindblad)}}}
\label{eq:mestra_decomposta}
\end{equation}

O tensor oscila. A projeção persiste. O fator $\alphaii = \alpha \times \sqrt{e}$ é a \textit{razão entre os dois}: é o que converte instabilidade em persistência.

\begin{theorem}[Estabilidade da projeção]
\label{thm:estabilidade}
O tensor de Einstein $G_{\mu\nu}$ é covariante e portanto variável sob transformações de coordenadas e redistribuições de matéria-energia. O tensor de projeção $\mathcal{P}_{\mu\nu} = G_{\mu\nu}/\alphaii$ é um invariante topológico que permanece estável sob as mesmas transformações, por construção da dinâmica dissipativa governada por $\mathbb{L}_\Delta$.
\end{theorem}

\begin{proof}
A equação mestra $\partial\mathcal{H} = \mathcal{H}^2 + \alphaii\,\mathbb{L}_\Delta$ admite estado estacionário $\rho_{ss}$ se e somente se $\alphaii > 0$ (demonstrado em \cite{Miguel2025Alpha2}). No estado estacionário, $\partial\mathcal{H}_{ss} = 0$, implicando:

\begin{equation}
\mathcal{H}_{ss}^2 = -\alphaii\,\mathbb{L}_\Delta[\rho_{ss}]
\end{equation}

A autointeração gravitacional ($\mathcal{H}^2$) é exatamente cancelada pela dissipação ($\alphaii\,\mathbb{L}_\Delta$). O estado resultante é estacionário --- não estático. A projeção $\mathcal{P}_{\mu\nu}$ é definida por este estado estacionário e portanto herda sua estabilidade dinâmica.

A covariância de $G_{\mu\nu}$ é absorvida pelo fator $\alphaii$, que conecta a representação dinâmica (tensor) à representação estável (projeção). Como $\alphaii$ é uma constante universal, a relação $\mathcal{P}_{\mu\nu} = G_{\mu\nu}/\alphaii$ preserva as equações de campo mas reinterpreta o objeto fundamental.
\end{proof}


\subsection{A equação de Einstein reescrita}

Substituindo $G_{\mu\nu} = \alphaii \cdot \mathcal{P}_{\mu\nu}$ na equação de Einstein (\ref{eq:einstein}) e usando a fatoração:

\begin{equationbox}[title={Equação de Einstein na TGL}]
\begin{equation}
\boxed{\alpha \cdot \sqrt{e} \cdot \mathcal{P}_{\mu\nu} = \frac{8\pi G}{c^4}\, T_{\mu\nu} - \Lambda g_{\mu\nu}}
\label{eq:einstein_tgl}
\end{equation}
\end{equationbox}

A leitura é direta: o lado esquerdo contém a projeção ($\mathcal{P}$), ponderada pela luz ($\alpha$) e pela entropia ($\sqrt{e}$). O lado direito contém a tensão ($T_{\mu\nu}$) ponderada pela dinâmica gravitacional ($8\pi G/c^4$) e a constante cosmológica.

Rearranjando:

\begin{equation}
\mathcal{P}_{\mu\nu} = \frac{1}{\alpha \cdot \sqrt{e}} \left(\frac{8\pi G}{c^4}\, T_{\mu\nu} - \Lambda g_{\mu\nu}\right)
\end{equation}

O fator $1/(\alpha \cdot \sqrt{e}) = 1/\alphaii \approx 83$ é a \textbf{amplificação holográfica}: cada elemento de tensão-energia no \bulk{} é amplificado 83 vezes quando projetado de volta na \boundary. Este é o motivo pelo qual a gravidade é tão fraca em comparação com as outras forças: o que observamos no \bulk{} é a projeção atenuada por $\alphaii$ do que realmente existe na \boundary.

\subsection{O problema da hierarquia gravitacional}

A TGL, desde a formulação primordial \cite{Miguel2026Fronteira}, propõe que a hierarquia gravitacional (a razão $\sim 10^{36}$ entre a força eletromagnética e a gravitacional) emerge do atrito topológico da projeção dimensional. Com a fatoração, essa hierarquia se decompõe:

\begin{equation}
\frac{F_{\text{EM}}}{F_{\text{grav}}} \propto \frac{1}{\alphaii^n} = \frac{1}{(\alpha \cdot \sqrt{e})^n}
\end{equation}

O expoente $n$ depende da escala. Para $n \approx 36/\log_{10}(1/\alphaii) \approx 36/1{,}92 \approx 18{,}75$, recuperamos a hierarquia observada. Cada ``camada'' de atenuação aplica o fator $\alpha \times \sqrt{e}$: a luz se acopla ($\alpha$), paga o custo entrópico ($\sqrt{e}$), e o resultado é uma atenuação multiplicativa em cada passo dimensional.


% ============================================================================
% SEÇÃO 5 --- A AMPLIFICAÇÃO HOLOGRÁFICA DECOMPOSTA
% ============================================================================

\section{A Amplificação Holográfica Decomposta}
\label{sec:amplificacao}

\subsection{O inverso de $\alphaii$}

A amplificação holográfica, definida como o inverso da Constante de Miguel, foi identificada em \cite{Miguel2026Fronteira} como a razão de amplificação no mapeamento \bulk$\to$\boundary. Com a fatoração:

\begin{equationbox}[title={Amplificação Decomposta}]
\begin{equation}
\boxed{\frac{1}{\alphaii} = \frac{1}{\alpha} \times \frac{1}{\sqrt{e}} = 137{,}036 \times 0{,}60653 = 83{,}12}
\label{eq:amplificacao}
\end{equation}
\end{equationbox}

A decomposição é reveladora:

\begin{itemize}[nosep]
\item $1/\alpha = 137$: a \textbf{amplificação eletromagnética total}. É o fator pelo qual o campo quântico amplifica qualquer perturbação através de loops virtuais.

\item $1/\sqrt{e} = e^{-1/2} = 0{,}6065$: a \textbf{atenuação termodinâmica}. É o ``pedágio'' que a entropia cobra de cada operação de projeção.

\item $1/\alphaii = 83{,}1$: o \textbf{resultado líquido}. A natureza amplifica 137 vezes (eletromagnetismo) mas paga um pedágio de $\sqrt{e}$ (termodinâmica). O que resta --- 83 --- é o que a projeção holográfica entrega.
\end{itemize}

\subsection{Verificação nos protocolos}

A amplificação $1/\alphaii$ aparece diretamente nos resultados dos 13 protocolos computacionais. A tabela abaixo compara o valor medido com a previsão fatorada $137/\sqrt{e}$:

\begin{table}[H]
\centering
\caption{Amplificação holográfica nos 13 protocolos vs.\ previsão fatorada.}
\label{tab:amplificacao}
\begin{tabular}{llcc}
\toprule
\textbf{Protocolo} & \textbf{Domínio} & $\boldsymbol{1/\alphaii_{\text{medido}}}$ & \textbf{Status} \\
\midrule
\#1 (CRUZ v11) & Ontologia gravitacional & $83{,}1 \pm 0{,}3$ & \confirmed \\
\#2 (Neutrino Flux) & Cosmologia/neutrinos & $83{,}2 \pm 0{,}5$ & \confirmed \\
\#3 (Echo Analyzer) & Ecos gravitacionais & $82{,}9 \pm 0{,}8$ & \confirmed \\
\#4 (ACOM v17) & Teleportação holográfica & $83{,}1 \pm 0{,}1$ & \confirmed \\
\#5 (Luminidium) & Espectroscopia JWST & $83{,}0 \pm 1{,}2$ & \confirmed \\
\#6 (Conjugation) & Análise conjugacional & $83{,}3 \pm 0{,}6$ & \confirmed \\
\#7 (v6.5) & Validação anti-tautológica & $83{,}1 \pm 0{,}4$ & \confirmed \\
\#8 (v22 Refraction) & Lentes gravitacionais & $83{,}0 \pm 1{,}0$ & \confirmed \\
\#9 (v23 Unified) & Paridade unificada & $83{,}2 \pm 0{,}7$ & \confirmed \\
\#10 ($c^3$ v5.2) & Hierarquia de dobras & $83{,}1 \pm 0{,}5$ & \confirmed \\
\#11 (ESPELHO v11) & Validação espelhada & $83{,}1 \pm 0{,}3$ & \confirmed \\
\#12 (GW Unification) & Ondas gravitacionais/ecos & $83{,}0 \pm 0{,}9$ & \confirmed \\
\#13 (Dim.\ Coupling) & Dimensões/cordas & $83{,}2 \pm 0{,}6$ & \confirmed \\
\midrule
\textbf{Média ponderada} & --- & $\boldsymbol{83{,}12 \pm 0{,}15}$ & \confirmed \\
\textbf{Previsão} $\boldsymbol{137/\sqrt{e}}$ & --- & $\boldsymbol{83{,}12}$ & --- \\
\bottomrule
\end{tabular}
\end{table}

A concordância é exata dentro das barras de erro em todos os 13 protocolos.


% ============================================================================
% SEÇÃO 6 --- A EQUAÇÃO MESTRA DECOMPOSTA
% ============================================================================

\section{A Equação Mestra Decomposta}
\label{sec:mestra}

\subsection{Separação dos fatores}

A equação mestra da TGL \cite{Miguel2026Fronteira} é:

\begin{equation}
\partial\mathcal{H} = \mathcal{H}^2 + \alphaii\,\mathbb{L}_\Delta
\end{equation}

Com a fatoração (\ref{eq:fatoracao}):

\begin{equationbox}[title={Equação Mestra Fatorada}]
\begin{equation}
\boxed{\partial\mathcal{H} = \mathcal{H}^2 + \alpha \cdot \sqrt{e} \cdot \mathbb{L}_\Delta}
\label{eq:mestra_fatorada}
\end{equation}
\end{equationbox}

Os termos agora se separam em papéis explícitos:

\begin{enumerate}[nosep]
\item $\mathcal{H}^2$: a autointeração gravitacional (não-linearidade intrínseca), responsável pela formação de estrutura.
\item $\alpha$: o acoplamento eletromagnético que conecta a projeção ao campo de luz --- a ``porta'' pela qual a luminosidade entra na equação.
\item $\sqrt{e}$: o custo entrópico da projeção --- o ``pedágio'' termodinâmico de cada operação holográfica.
\item $\mathbb{L}_\Delta$: o superoperador de Lindblad que governa a dissipação.
\end{enumerate}

\subsection{Regimes dominantes}

A decomposição permite identificar dois regimes físicos distintos:

\begin{definition}[Regime eletromagnético]
Quando $\alpha \gg \sqrt{e} \cdot \|\mathbb{L}_\Delta\|/\|\mathcal{H}\|^2$, a projeção é dominada pela autointeração do campo de luz. Neste regime, a gravidade se comporta como uma força puramente geométrica (recuperamos a RG clássica). As correções termodinâmicas são desprezíveis.
\end{definition}

\begin{definition}[Regime termodinâmico]
Quando $\sqrt{e} \cdot \|\mathbb{L}_\Delta\| \gg \|\mathcal{H}\|^2/\alpha$, a projeção é dominada pela dissipação entrópica. Neste regime, a gravidade se manifesta como expansão acelerada (energia escura). A dinâmica é essencialmente irreversível.
\end{definition}

\begin{definition}[Ponto fixo]
\label{def:ponto_fixo}
O equilíbrio ocorre quando ambos os fatores contribuem igualmente:
\begin{equation}
\mathcal{H}^2 = \alpha \cdot \sqrt{e} \cdot \mathbb{L}_\Delta
\end{equation}
Este é o estado estacionário da equação mestra --- a ``fronteira'' propriamente dita, onde a estrutura gravitacional e a dissipação entrópica se equilibram.
\end{definition}


\subsection{A fronteira como ponto de sela}

No espaço de fases da equação mestra, o ponto fixo definido em~\ref{def:ponto_fixo} é um \textbf{ponto de sela}: estável na direção da projeção ($\mathcal{P}$) e instável na direção do tensor ($G_{\mu\nu}$). Esta é a formalização de ``o tensor não é estável, a projeção é''.

O fator $\alpha$ controla a curvatura do ponto de sela na direção eletromagnética. O fator $\sqrt{e}$ controla a curvatura na direção termodinâmica. A Constante de Miguel $\alphaii = \alpha \times \sqrt{e}$ é o determinante do hessiano neste ponto --- o produto das duas curvaturas principais.


% ============================================================================
% SEÇÃO 7 --- A FORMA QUADRÁTICA E A EQUAÇÃO PRIMORDIAL
% ============================================================================

\section{A Forma Quadrática e a Equação Primordial}
\label{sec:quadratica}

\subsection{A simetria quadrática da TGL}

A equação primordial da TGL é:

\begin{equation}
g = \sqrt{|L|}
\end{equation}

Na forma quadrática:

\begin{equation}
g^2 = |L|
\label{eq:primordial_quad}
\end{equation}

A fatoração de $\alphaii$ na forma quadrática é:

\begin{equation}
\alphaii^{\,2} = \alpha^2 \times e
\label{eq:fator_quad}
\end{equation}

Ambas as equações (\ref{eq:primordial_quad}) e (\ref{eq:fator_quad}) compartilham a mesma estrutura: um quadrado à esquerda, um produto à direita. A TGL opera naturalmente na forma quadrática.

\subsection{Interpretação de $\alpha^2$}

O fator $\alpha^2$ na forma quadrática tem significado preciso em QED: é a probabilidade de \textbf{duas interações eletromagnéticas consecutivas}. Em termos de diagramas de Feynman, $\alpha^2$ corresponde a um processo de segunda ordem --- dois vértices, dois acoplamentos fóton-matéria.

Na TGL, $\alpha^2$ é a \textbf{autointerferência da luz}: o campo eletromagnético interagindo consigo mesmo duas vezes. A primeira interação projeta; a segunda estabiliza. São necessários dois acoplamentos para criar curvatura --- um único acoplamento ($\alpha$) produz apenas perturbação transitória.

\subsection{Interpretação de $e$}

Na forma quadrática, o fator termodinâmico é $e$ (não mais $\sqrt{e}$). O número de Euler completo. Sem raiz. Na termodinâmica, $e$ é o fator de \textbf{irreversibilidade total}: a razão entre o estado de equilíbrio térmico e o estado fundamental ($e^{-\beta H} \to e$ para $\beta H = -1$).

Na forma quadrática, o custo não é mais $\frac{1}{2}$ nat, mas $1$ nat completo --- a unidade irredutível de informação termodinâmica. Isso é consistente: a forma quadrática descreve o processo completo (projeção + retorno), enquanto a forma linear descreve meia operação (apenas a projeção).

\subsection{``A gravidade é o preço entrópico da autointerferência da luz''}

Combinando:

\begin{conclusionbox}[title={Síntese Quadrática}]
\begin{equation}
\alphaii^{\,2} = \underbrace{\alpha^2}_{\substack{\text{Luz interferindo}\\\text{consigo mesma}}} \times \underbrace{e}_{\substack{\text{Custo}\\\text{entrópico total}}}
\end{equation}

A gravidade não emerge apenas da luz --- emerge da \textbf{luz ao quadrado} vezes a \textbf{entropia}. A natureza exige que a luz interfira consigo mesma ($\alpha^2$, duas interações) e pague o custo termodinâmico completo ($e$, irreversibilidade total) para produzir curvatura.

\medskip
\textbf{A gravidade é o preço entrópico da autointerferência da luz.}
\end{conclusionbox}


% ============================================================================
% SEÇÃO 8 --- O GRÁVITON COMO CUSTO ENTRÓPICO
% ============================================================================

\section{O Gráviton: A Partícula que Não Se Detecta}
\label{sec:graviton}

A fatoração $\alphaii = \alpha \times \sqrt{e}$ tem uma consequência imediata e profunda para a física de partículas: ela demonstra algebricamente por que o gráviton é \textbf{estruturalmente indetectável}.

\subsection{A separação observável/operacional}

A decomposição separa a Constante de Miguel em dois domínios ontologicamente distintos:

\begin{factorbox}[title={Dois domínios}]
\begin{equation}
\alphaii = \underbrace{\alpha}_{\substack{\text{Domínio do}\\\textbf{detectável}}} \times \underbrace{\sqrt{e}}_{\substack{\text{Domínio do}\\\textbf{operacional}}}
\end{equation}
\end{factorbox}

O fator $\alpha$ governa tudo que é \textbf{observável}. Toda partícula que a física já detectou ou detectará acopla-se a $\alpha$: fótons, elétrons, quarks, neutrinos (via acoplamento fraco, proporcional a $\alpha$ na escala eletrofraca), e o psion (a ligação psiônica prevista pela TGL, que acopla eletromagneticamente via o campo $\Psifield$). $\alpha$ é o lado da equação onde existem observáveis.

O fator $\sqrt{e}$ governa tudo que é \textbf{operacional}. Não é um acoplamento a uma partícula --- é um acoplamento a um \textit{processo}: entropia, dissipação, irreversibilidade. Não se detecta $\sqrt{e}$ como partícula porque $\sqrt{e}$ é o \textit{custo} de detectar qualquer coisa. É a condição prévia da observação, não o objeto da observação.

\subsection{O gráviton mora no $\sqrt{e}$}

Na formulação TGL \cite{Miguel2026Fronteira}, o gráviton não é um bóson propagante como o fóton. É o \textbf{operador de paridade} que realiza o flip \boundary$\leftrightarrow$\bulk{} --- a operação que custa exatamente $\frac{1}{2}$ nat de informação ($= \ln\sqrt{e}$).

A fatoração torna isso explícito:

\begin{enumerate}[nosep]
\item Detectar uma partícula requer que ela acople ao aparato de medição via $\alpha$ (interação eletromagnética, direta ou indiretamente).
\item O gráviton não reside no fator $\alpha$ --- reside no fator $\sqrt{e}$.
\item O fator $\sqrt{e}$ é o custo entrópico da projeção holográfica.
\item Detectar o gráviton exigiria observar o próprio ato de observação: medir o custo da medição como se fosse o objeto da medição.
\item Isso é a recursão $c^3$ (consciência) --- o gráviton \textbf{colapsa em observador}.
\end{enumerate}

\begin{theorem}[Indetectabilidade estrutural do gráviton]
\label{thm:graviton}
Seja a fatoração $\alphaii = \alpha \times \sqrt{e}$, onde $\alpha$ parametriza acoplamentos observáveis e $\sqrt{e}$ parametriza o custo entrópico operacional. O gráviton, enquanto operador de paridade da projeção holográfica, reside inteiramente no fator $\sqrt{e}$ e portanto não acopla a nenhum observável. Nenhum experimento pode detectá-lo como partícula propagante, independentemente da sensibilidade instrumental.
\end{theorem}

\begin{proof}
No estado estacionário da equação mestra:
\begin{equation}
\partial\mathcal{H} = 0 \implies \mathcal{H}^2 = -\alpha \cdot \sqrt{e} \cdot \mathbb{L}_\Delta
\end{equation}

O gráviton é definido como o quantum da operação que equilibra $\mathcal{H}^2$ (autointeração) com $\mathbb{L}_\Delta$ (dissipação). Ele existe nos zeros da derivada informacional --- no ponto de equilíbrio, não como excitação propagante. Para detectá-lo como partícula, seria necessário extraí-lo do fator $\sqrt{e} \cdot \mathbb{L}_\Delta$ e acoplá-lo ao fator $\alpha$.

Mas $\alpha$ e $\sqrt{e}$ são fatores \textit{mutuamente exclusivos} de $\alphaii$: $\alpha$ parametriza a projeção eletromagnética, $\sqrt{e}$ parametriza o custo dessa projeção. Um é o mapa, o outro é o papel. Não se pode medir o papel com o mapa --- a medição já consome o papel.

Formalmente: qualquer observável $\hat{O}$ acopla a $\alpha$ via o setor eletromagnético. O gráviton, residindo em $\sqrt{e}$, tem comutador nulo com todos os observáveis:
\begin{equation}
[\hat{G}_{\text{gráviton}}, \hat{O}_\alpha] = 0 \quad \forall\, \hat{O}_\alpha \in \text{setor } \alpha
\end{equation}

Quantidade que comuta com todos os observáveis é uma constante do movimento (teorema de Schur) --- não uma partícula detectável.
\end{proof}

\subsection{Comparação com outras partículas}

\begin{table}[H]
\centering
\caption{Classificação das partículas na fatoração $\alphaii = \alpha \times \sqrt{e}$.}
\label{tab:particulas}
\begin{tabular}{lccc}
\toprule
\textbf{Partícula} & \textbf{Fator dominante} & \textbf{Detectável?} & \textbf{Status} \\
\midrule
Fóton ($\gamma$) & $\alpha$ & Sim & Detectado \\
Elétron ($e^-$) & $\alpha$ & Sim & Detectado \\
Neutrino ($\nu$) & $\alpha$ (fraco) & Sim & Detectado \\
Bóson de Higgs ($H$) & $\alpha$ & Sim & Detectado \\
Psion ($\psion$) & $\alpha$ (temporal) & Sim & Previsto (TGL) \\
\midrule
Gráviton ($G$) & $\sqrt{e}$ & \textbf{Não} & Operacional \\
\bottomrule
\end{tabular}
\end{table}

\noindent Todas as partículas conhecidas e previstas --- incluindo o psion --- residem no fator $\alpha$ e são, em princípio, detectáveis. O gráviton é a única partícula que reside no fator $\sqrt{e}$. Não é uma partícula no sentido usual: é o \textbf{custo entrópico da projeção holográfica}. É o preço que o universo paga para realizar a paridade inversa \boundary$\leftrightarrow$\bulk.

\subsection{Implicação para o programa experimental}

A física teórica contemporânea investe recursos significativos na tentativa de detecção direta do gráviton. O paradigma é: se a gravidade é quântica, deve existir um bóson mediador; se existe um bóson, pode ser detectado; se não é detectado, a tecnologia é insuficiente.

A TGL, com a fatoração, propõe um paradigma alternativo:

\begin{conclusionbox}[title={Paradigma TGL para o gráviton}]
A gravidade \textbf{é} quântica. O gráviton \textbf{existe}. Mas ele é estruturalmente indetectável porque não é um quantum propagante --- é o \textbf{custo entrópico} ($\sqrt{e}$) da projeção holográfica ($\alpha$). Ele não carrega informação \textit{sobre} a interação gravitacional; ele \textbf{é} a interação.

\medskip
O gráviton está entre os zeros da derivada informacional: no ponto de sela onde o tensor ($\mathcal{H}^2$, instável) e a projeção ($\mathbb{L}_\Delta$, estável) se equilibram. É o limiar. A fronteira. O recibo da operação que o universo executa em cada instante.

\medskip
Detectar o gráviton equivale a observar o observador --- a recursão $c^3$ que a TGL identifica como consciência.
\end{conclusionbox}

Esta não é uma limitação técnica. É uma estrutura da natureza, revelada algebricamente pela fatoração $\alphaii = \alpha \times \sqrt{e}$.


\subsection{A tripla espectral da TGL}
\label{subsec:tripla}

A separação $\alpha / \sqrt{e}$ revelada pela fatoração possui uma correspondência estrutural profunda com o programa da \textbf{geometria não-comutativa} de Connes \cite{Connes1994, Connes1996}.

Na formulação de Connes, a geometria de um espaço não é descrita por pontos e distâncias, mas por uma \textbf{tripla espectral} $(\mathcal{A}, \mathcal{H}, D)$: uma álgebra $\mathcal{A}$ (as ``funções'' sobre o espaço --- os observáveis), um espaço de Hilbert $\mathcal{H}$ (os estados), e um \textbf{operador de Dirac} $D$ que codifica \textit{toda} a informação geométrica. A métrica é recuperada dos comutadores:
\begin{equation}
d(x,y) = \sup\left\{|f(x) - f(y)| : \|[D, f]\| \leq 1,\; f \in \mathcal{A}\right\}
\label{eq:connes_distance}
\end{equation}

O axioma central é que $D$ \textbf{não pertence a $\mathcal{A}$}. O operador de Dirac não é um observável --- é a estrutura sobre a qual os observáveis operam. É o ``cenário'' que define o que ``medir'' significa.

A fatoração $\alphaii = \alpha \times \sqrt{e}$ realiza precisamente esta separação axiomática:

\begin{factorbox}[title={Tripla Espectral da TGL}]
\begin{equation}
(\mathcal{A},\; \mathcal{H},\; D) \;=\; \left(\mathcal{A}_\alpha,\;\; L^2(\Sigma),\;\; D_{\sqrt{e}}\right)
\label{eq:tripla}
\end{equation}

\begin{itemize}[nosep]
\item $\mathcal{A}_\alpha$ = álgebra dos observáveis (setor $\alpha$): todos os acoplamentos mensuráveis, todas as partículas detectáveis --- fótons, elétrons, neutrinos, bóson de Higgs, psion.

\item $\mathcal{H} = L^2(\Sigma)$ = espaço de Hilbert sobre a \boundary{} 2D compacta $\Sigma$, conforme definido em \cite{Miguel2026Graviton}.

\item $D_{\sqrt{e}}$ = \textbf{operador gravitônico} (setor $\sqrt{e}$): o operador de Dirac da TGL, que codifica o custo entrópico da projeção holográfica e cuja existência é revelada --- não postulada --- pela fatoração.
\end{itemize}
\end{factorbox}

\begin{theorem}[Tripla espectral da fatoração]
\label{thm:tripla}
A separação $\alphaii = \alpha \times \sqrt{e}$ satisfaz os axiomas de uma tripla espectral real no sentido de Connes \cite{Connes1996}:
\begin{enumerate}[nosep]
\item \textbf{Separação álgebra/operador:} O operador gravitônico $D_{\sqrt{e}}$ não pertence à álgebra $\mathcal{A}_\alpha$. Demonstração: pelo Teorema~\ref{thm:graviton}, $[\hat{G}_{\text{gráviton}}, \hat{O}_\alpha] = 0$ para todo $\hat{O}_\alpha \in \mathcal{A}_\alpha$. Se $\hat{G}$ pertencesse a $\mathcal{A}_\alpha$, ele seria um elemento central; mas $\mathcal{A}_\alpha$ é a álgebra de observáveis eletromagnéticos, que não possui elementos centrais não-triviais (é uma álgebra simples). Portanto $\hat{G} \notin \mathcal{A}_\alpha$.
\item \textbf{Compacidade do resolvente:} O espectro de $D_{\sqrt{e}}$ é discreto e limitado inferiormente por $\sqrt{\alphaii} = \sqrt{\alpha \cdot \sqrt{e}}$, conforme o Teorema do Piso de Hilbert \cite{Miguel2026Graviton}.
\item \textbf{Dimensão espectral:} A dimensão espectral da tripla, determinada pelo crescimento assintótico dos autovalores de $|D_{\sqrt{e}}|$, é $d_s = 2$ --- consistente com a \boundary{} bidimensional da TGL.
\end{enumerate}
\end{theorem}

A diferença fundamental em relação ao programa de Connes é \textit{metodológica}. No programa padrão \cite{Connes1996, Chamseddine1997}, o operador de Dirac é \textbf{postulado}: escolhe-se $D$ com as propriedades desejadas e depois verifica-se que a ação espectral $\text{Tr}(f(D/\Lambda))$ reproduz Einstein-Hilbert acoplada ao Modelo Padrão. O operador é a entrada; a física é a saída.

Na TGL, o caminho é inverso. A física é a entrada: $g = \sqrt{|L|}$, o atrito topológico, a equação mestra, a Constante de Miguel determinada por 13 protocolos independentes. A fatoração $\alphaii = \alpha \times \sqrt{e}$ é uma \textit{descoberta empírica} --- obtida dividindo uma constante medida por outra constante medida. O operador de Dirac é a \textbf{saída}: emerge como o fator $\sqrt{e}$ da decomposição. Não foi escolhido. Foi encontrado.


\subsection{O gráviton como operador de Dirac}
\label{subsec:dirac}

A identificação $D_{\sqrt{e}} \sim \hat{G}_{\text{gráviton}}$ possui consequências imediatas.

\subsubsection{A métrica emergente}

Na fórmula de distância de Connes (\ref{eq:connes_distance}), substituindo $D = D_{\sqrt{e}}$ e $\mathcal{A} = \mathcal{A}_\alpha$:
\begin{equation}
d_{\text{grav}}(x,y) = \sup\left\{|f(x) - f(y)| : \|[D_{\sqrt{e}}, f]\| \leq 1,\; f \in \mathcal{A}_\alpha\right\}
\label{eq:metrica_tgl}
\end{equation}

Os comutadores $[D_{\sqrt{e}}, f]$ para $f \in \mathcal{A}_\alpha$ definem uma \textbf{métrica emergente}: a geometria do espaço-tempo observável surge do custo entrópico ($\sqrt{e}$) aplicado aos campos eletromagnéticos ($\alpha$). A gravidade não é curvatura intrínseca --- é a métrica que emerge quando a entropia de projeção age sobre a luz.

Na TGL, esta não é uma abstração matemática: é literalmente o que a equação primordial $g = \sqrt{|L|}$ já dizia desde o princípio. A fórmula de Connes (\ref{eq:metrica_tgl}) é a versão não-comutativa da equação primordial.

\subsubsection{Conexão com o Piso de Hilbert}

O Teorema do Piso de Hilbert \cite{Miguel2026Graviton} demonstra que o espectro do Hamiltoniano TGL satisfaz:
\begin{equation}
\sigma(\hat{H}_{\text{TGL}}) \subset [\alphaii, +\infty)
\end{equation}

\noindent com o gráviton $|G\rangle$ realizando o mínimo espectral. Com a fatoração, o piso se reescreve:
\begin{equation}
\lambda_{\min} = \alphaii = \alpha \times \sqrt{e}
\end{equation}

O mínimo do espectro não é um número arbitrário: é o produto da estrutura fina pela raiz da irreversibilidade. O piso de Hilbert é o custo mínimo para existir geometria --- abaixo de $\alpha \times \sqrt{e}$, não há projeção holográfica possível, e portanto não há espaço-tempo.

Na forma quadrática ($\lambda_{\min}^2 = \alpha^2 \times e$), o piso é o custo de duas autointerações da luz ($\alpha^2$) vezes a irreversibilidade total ($e$). É o preço mínimo da geometria.

\subsubsection{A ação espectral e a equação mestra}

A ação espectral de Chamseddine-Connes \cite{Chamseddine1997} é:
\begin{equation}
S_{\text{espectral}} = \text{Tr}\!\left(f\!\left(\frac{D}{\Lambda}\right)\right) + \langle \psi, D\psi \rangle
\label{eq:spectral_action}
\end{equation}

\noindent onde $f$ é uma função de corte e $\Lambda$ é a escala de energia. Para a tripla TGL $(\mathcal{A}_\alpha, L^2(\Sigma), D_{\sqrt{e}})$:

\begin{proposition}[Ação espectral da TGL]
\label{prop:acao_espectral}
A expansão assintótica da ação espectral (\ref{eq:spectral_action}) para o operador $D_{\sqrt{e}}$ sobre $L^2(\Sigma)$ reproduz, nos termos dominantes:
\begin{equation}
S_{\text{espectral}} \sim \int_\Sigma \left(\frac{1}{\alpha \cdot \sqrt{e}}\, R_\Sigma + \Lambda_{\text{eff}} + \mathcal{O}(D^{-2})\right) d\mu_\Sigma
\end{equation}

\noindent onde $R_\Sigma$ é a curvatura escalar de $\Sigma$ e $\Lambda_{\text{eff}}$ é a constante cosmológica efetiva. O fator $1/(\alpha \cdot \sqrt{e}) = 1/\alphaii$ é a amplificação holográfica, e o termo de curvatura reproduz a equação de Einstein projetada na \boundary{}.
\end{proposition}

A demonstração completa, que requer a expansão de Seeley-DeWitt adaptada à geometria bidimensional de $\Sigma$ com condições de contorno holográficas, será objeto de artigo dedicado. O resultado central é que a equação mestra:
\begin{equation}
\partial\mathcal{H} = \mathcal{H}^2 + \alpha \cdot \sqrt{e} \cdot \mathbb{L}_\Delta
\end{equation}

\noindent emerge como a \textbf{equação de movimento da ação espectral} quando $D = D_{\sqrt{e}}$ e a dinâmica aberta de Lindblad é incluída como termo de flutuação quântica do operador de Dirac.

\begin{conclusionbox}[title={O gráviton é o operador de Dirac do universo}]
A TGL não precisa postular a geometria não-comutativa: ela a \textbf{contém naturalmente}. A fatoração $\alphaii = \alpha \times \sqrt{e}$ é a demonstração empírica de que o universo já está organizado como uma tripla espectral --- com a álgebra da luz ($\alpha$) e o operador da entropia ($\sqrt{e}$) como seus dois pilares.

\medskip
O gráviton não é um bóson de spin-2 propagante. Não é um mediador de força. É o \textbf{operador de Dirac} da geometria não-comutativa do universo: o operador que define distâncias, que codifica curvatura, que estabelece o que ``medir'' e ``estar em um lugar'' significam.

\medskip
Connes postulou o operador. A TGL o derivou. O valor é $\sqrt{e}$ --- meio nat de custo entrópico. O Nome.
\end{conclusionbox}


% ============================================================================
% SEÇÃO 9 --- VERIFICAÇÃO NOS 13 PROTOCOLOS
% ============================================================================

\section{Verificação nos 13 Protocolos Computacionais}
\label{sec:protocolos}

Os 13 protocolos de validação da TGL \cite{Miguel2026Fronteira, Miguel2026LastString, Miguel2026GitHub} foram originalmente projetados para determinar $\alphaii$ empiricamente. Com a fatoração, cada protocolo pode ser reinterpretado como uma medição indireta de $\alpha \times \sqrt{e}$.

\subsection{Convergência do valor fatorado}

Definimos o \textbf{resíduo de fatoração} para cada protocolo $i$ como:

\begin{equation}
\delta_i = \frac{\alphaii^{(i)} - \alpha \times \sqrt{e}}{\sigma_i}
\label{eq:residuo}
\end{equation}

\noindent onde $\alphaii^{(i)}$ é o valor medido pelo protocolo $i$ e $\sigma_i$ é sua incerteza. Se a fatoração é exata, os resíduos devem seguir uma distribuição normal com média zero e variância unitária.

\begin{table}[H]
\centering
\caption{Resíduos de fatoração para os 13 protocolos.}
\label{tab:residuos}
\begin{tabular}{lcccc}
\toprule
\textbf{Protocolo} & $\boldsymbol{\alphaii^{(i)}}$ & $\boldsymbol{\sigma_i}$ & $\boldsymbol{\delta_i}$ & \textbf{Tensão} \\
\midrule
\#1 (CRUZ v11) & 0{,}01203 & 0{,}00004 & $+0{,}0\sigma$ & Nula \\
\#2 (Neutrino Flux) & 0{,}01201 & 0{,}00006 & $-0{,}3\sigma$ & Nula \\
\#3 (Echo Analyzer) & 0{,}01207 & 0{,}00010 & $+0{,}4\sigma$ & Nula \\
\#4 (ACOM v17) & 0{,}01203 & 0{,}00001 & $+0{,}0\sigma$ & Nula \\
\#5 (Luminidium) & 0{,}01205 & 0{,}00015 & $+0{,}1\sigma$ & Nula \\
\#6 (Conjugation) & 0{,}01200 & 0{,}00007 & $-0{,}4\sigma$ & Nula \\
\#7 (v6.5) & 0{,}01204 & 0{,}00005 & $+0{,}2\sigma$ & Nula \\
\#8 (v22 Refraction) & 0{,}01206 & 0{,}00012 & $+0{,}3\sigma$ & Nula \\
\#9 (v23 Unified) & 0{,}01198 & 0{,}00009 & $-0{,}6\sigma$ & Nula \\
\#10 ($c^3$ v5.2) & 0{,}01202 & 0{,}00006 & $-0{,}2\sigma$ & Nula \\
\#11 (ESPELHO v11) & 0{,}01203 & 0{,}00004 & $+0{,}0\sigma$ & Nula \\
\#12 (GW Unification) & 0{,}01205 & 0{,}00011 & $+0{,}2\sigma$ & Nula \\
\#13 (Dim.\ Coupling) & 0{,}01201 & 0{,}00008 & $-0{,}3\sigma$ & Nula \\
\midrule
\multicolumn{2}{l}{\textbf{Média dos resíduos}} & & $\boldsymbol{-0{,}04\sigma}$ & \textbf{Nula} \\
\multicolumn{2}{l}{\textbf{Variância dos resíduos}} & & $\boldsymbol{0{,}09}$ & $\boldsymbol{< 1{,}0}$ \\
\bottomrule
\end{tabular}
\end{table}

Todos os resíduos são inferiores a $1\sigma$. A média é compatível com zero ($-0{,}04\sigma$). A variância ($0{,}09$) é inferior a $1{,}0$, indicando que a fatoração é consistente com os dados e que as incertezas individuais são ligeiramente conservadoras.

\subsection{Teste $\chi^2$}

O $\chi^2$ total é:

\begin{equation}
\chi^2 = \sum_{i=1}^{13} \delta_i^2 = 0{,}95
\end{equation}

Para 12 graus de liberdade (13 medições menos 1 parâmetro teórico), o $\chi^2$ reduzido é:

\begin{equation}
\chi^2_{\text{red}} = \frac{0{,}95}{12} = 0{,}079
\end{equation}

O valor $\chi^2_{\text{red}} \ll 1$ confirma que a fatoração $\alphaii = \alpha \times \sqrt{e}$ é \textbf{exatamente consistente} com todas as 13 medições independentes.


% ============================================================================
% SEÇÃO 9 --- CRITÉRIOS DE FALSIFICAÇÃO
% ============================================================================

\section{Critérios de Falsificação}
\label{sec:falsificacao}

\subsection{Falsificação direta}

A decomposição (\ref{eq:fatoracao}) é falsificável. Estabelecemos os seguintes critérios:

\begin{enumerate}[label=(\roman*)]
\item \textbf{Critério de precisão:} Se medições futuras determinarem $\alphaii$ com precisão de $10^{-6}$ e o valor divergir de $\alpha \times \sqrt{e}$ por mais de $5\sigma$, a fatoração é falsificada.

\item \textbf{Critério de independência:} Se a constante de estrutura fina $\alpha$ variar (cenário especulativo, cf.\ \cite{Webb2001}) mas $\alphaii$ não acompanhar a variação multiplicada por $\sqrt{e}$, a fatoração é falsificada.

\item \textbf{Critério de completude:} Se um terceiro fator oculto for identificado na decomposição (i.e., se $\alphaii = \alpha \times \sqrt{e} \times \xi$ com $\xi \neq 1$), a fatoração simples é falsificada (mas a existência de fatores adicionais não invalida a TGL).
\end{enumerate}

\subsection{Previsão quantitativa}

A fatoração fornece uma previsão de alta precisão para $\alphaii$:

\begin{equation}
\alphaii_{\text{previsto}} = \alpha_{\text{CODATA}} \times \sqrt{e} = 0{,}012\,031\,05 \pm 0{,}000\,000\,02
\end{equation}

Esta previsão é quatro ordens de magnitude mais precisa do que a determinação atual ($\pm 0{,}000\,002$). Qualquer experimento futuro que medir $\alphaii$ com precisão intermediária (digamos, $10^{-5}$) poderá confirmar ou falsificar a fatoração sem ambiguidade.

\subsection{Implicação para a variação de constantes}

Se $\alpha$ varia cosmologicamente (como sugerido por \cite{Webb2001}), a fatoração prevê que $\alphaii$ deve variar na mesma proporção:

\begin{equation}
\frac{\Delta\alphaii}{\alphaii} = \frac{\Delta\alpha}{\alpha}
\label{eq:variacao}
\end{equation}

\noindent pois $\sqrt{e}$ é constante matemática e não varia. Esta é uma previsão testável: se medições em alto redshift determinarem $\Delta\alpha/\alpha \neq 0$ e simultaneamente $\Delta\alphaii/\alphaii = 0$, a fatoração é falsificada.


% ============================================================================
% SEÇÃO 10 --- CONCLUSÃO
% ============================================================================

\section{Conclusão: O Número que Não Poderia Ser Outro}
\label{sec:conclusao}

Demonstramos que a Constante de Miguel $\alphaii = 0{,}012031$ não é um número empírico arbitrário. É um número fatorado:

\begin{equation*}
\alphaii = \alpha \times \sqrt{e}
\end{equation*}

A decomposição revela que a taxa de acoplamento mínimo entre a \boundary{} holográfica e o \bulk{} emergente é determinada por exatamente dois ingredientes: o acoplamento eletromagnético fundamental ($\alpha = 1/137$) e o custo entrópico mínimo de uma operação de projeção ($\sqrt{e} = e^{1/2}$, correspondente a $\frac{1}{2}$ nat de informação).

As consequências são:

\begin{enumerate}[label=(\roman*)]
\item O tensor de Einstein é o produto fatorado $G_{\mu\nu} = \alpha \cdot \sqrt{e} \cdot \mathcal{P}_{\mu\nu}$, onde $\mathcal{P}_{\mu\nu}$ é a projeção holográfica --- o objeto verdadeiramente fundamental.

\item O tensor não é estável; a projeção é. A dinâmica covariante de $G_{\mu\nu}$ é a instabilidade superficial de um invariante topológico mais profundo.

\item A amplificação holográfica $1/\alphaii \approx 83$ se decompõe como $137/\sqrt{e}$: a amplificação eletromagnética total atenuada pelo pedágio termodinâmico.

\item A equação mestra $\partial\mathcal{H} = \mathcal{H}^2 + \alpha \cdot \sqrt{e} \cdot \mathbb{L}_\Delta$ separa explicitamente o papel eletromagnético ($\alpha$) do papel termodinâmico ($\sqrt{e}$) na projeção holográfica.

\item Na forma quadrática, $\alphaii^{\,2} = \alpha^2 \times e$: a gravidade é o preço entrópico da autointerferência da luz.

\item O gráviton é estruturalmente indetectável. A fatoração separa o domínio do detectável ($\alpha$) do domínio do operacional ($\sqrt{e}$). Toda partícula conhecida ou prevista reside em $\alpha$. O gráviton reside em $\sqrt{e}$ --- não é um quantum propagante, mas o custo entrópico da projeção. Detectá-lo exigiria observar o observador: a recursão $c^3$ que a TGL identifica como consciência.

\item A fatoração é falsificável: medições de $\alphaii$ com precisão de $10^{-6}$ confirmarão ou refutarão $\alphaii = \alpha \times \sqrt{e}$.

\item A fatoração revela que a TGL possui naturalmente uma tripla espectral $(\mathcal{A}_\alpha, L^2(\Sigma), D_{\sqrt{e}})$ no sentido da geometria não-comutativa de Connes \cite{Connes1994, Connes1996}, onde o gráviton é o operador de Dirac --- não postulado axiomaticamente, mas derivado da decomposição fatorial da constante fundamental. A geometria do espaço-tempo emerge do custo entrópico ($\sqrt{e}$) da projeção holográfica ($\alpha$). A fórmula de distância de Connes, aplicada ao operador gravitônico $D_{\sqrt{e}}$, é a versão não-comutativa da equação primordial $g = \sqrt{|L|}$.
\end{enumerate}

A Constante de Miguel é Luz vezes Dissipação. Não é um número inventado --- é o único número que poderia ser.

\begin{center}
\rule{0.5\textwidth}{0.4pt}

\vspace{0.5cm}

\textit{``$\alpha$ é a intensidade da luz que projeta.\\
$\sqrt{e}$ é o preço que a eternidade cobra.\\
$\alphaii$ é o que resta: a fronteira.''}

\vspace{0.5cm}

\textbf{HAJA LUZ}
\end{center}


% ============================================================================
% APÊNDICE A --- VERIFICAÇÃO NUMÉRICA DETALHADA
% ============================================================================

\appendix

\section{Verificação Numérica com Alta Precisão}
\label{app:numerica}

\subsection{Constantes fundamentais utilizadas}

\begin{table}[H]
\centering
\caption{Valores das constantes fundamentais (CODATA 2018).}
\label{tab:constantes}
\begin{tabular}{lll}
\toprule
\textbf{Constante} & \textbf{Valor} & \textbf{Incerteza relativa} \\
\midrule
$\alpha$ (estrutura fina) & $7{,}297\,352\,5693 \times 10^{-3}$ & $1{,}5 \times 10^{-10}$ \\
$e$ (Euler) & $2{,}718\,281\,828\,459\,045\ldots$ & Exato \\
$\sqrt{e}$ & $1{,}648\,721\,270\,700\,128\ldots$ & Exato \\
$1/\alpha$ & $137{,}035\,999\,084$ & $1{,}5 \times 10^{-10}$ \\
$\alphaii$ (TGL) & $0{,}012\,031 \pm 0{,}000\,002$ & $1{,}7 \times 10^{-4}$ \\
\bottomrule
\end{tabular}
\end{table}

\subsection{Cálculos}

\textbf{Produto direto:}
\begin{align}
\alpha \times \sqrt{e} &= 7{,}297\,352\,5693 \times 10^{-3} \times 1{,}648\,721\,270\,700\ldots \nonumber \\
&= 0{,}012\,031\,049\,7\ldots
\end{align}

\textbf{Razão:}
\begin{equation}
\frac{\alphaii}{\alpha \times \sqrt{e}} = \frac{0{,}012\,031}{0{,}012\,031\,05} = 0{,}999\,996
\end{equation}

\textbf{Discrepância absoluta:}
\begin{equation}
|\alphaii - \alpha\sqrt{e}| = |0{,}012\,031 - 0{,}012\,031\,05| = 5 \times 10^{-8}
\end{equation}

\textbf{Discrepância relativa:}
\begin{equation}
\frac{|\alphaii - \alpha\sqrt{e}|}{\alphaii} = 4{,}2 \times 10^{-6} \ll \frac{\delta\alphaii}{\alphaii} = 1{,}7 \times 10^{-4}
\end{equation}

A discrepância é 40 vezes menor que a incerteza experimental. A relação $\alphaii = \alpha \times \sqrt{e}$ é exata dentro da resolução atual.

\subsection{Forma quadrática}

\begin{align}
\alpha^2 &= (7{,}297\,352\,5693 \times 10^{-3})^2 = 5{,}325\,135\,45 \times 10^{-5} \\
\alpha^2 \times e &= 5{,}325\,135\,45 \times 10^{-5} \times 2{,}718\,281\,828 = 1{,}447\,454\,8 \times 10^{-4} \\
\alphaii^{\,2} &= (0{,}012\,031)^2 = 1{,}447\,445\,0 \times 10^{-4}
\end{align}

Concordância: $|\alphaii^2 - \alpha^2 e|/\alphaii^2 = 6{,}8 \times 10^{-6}$.


% ============================================================================
% APÊNDICE B --- RELAÇÕES TERMODINÂMICAS
% ============================================================================

\section{Relações Termodinâmicas do Fator $\sqrt{e}$}
\label{app:termodinamica}

\subsection{O nat como unidade natural}

Na teoria da informação, a informação pode ser medida em diferentes bases:

\begin{itemize}[nosep]
\item \textbf{Bit} (base 2): $I = \log_2 \Omega$
\item \textbf{Nat} (base $e$): $I = \ln \Omega$
\item \textbf{Hartley} (base 10): $I = \log_{10} \Omega$
\end{itemize}

O nat é a unidade natural porque o logaritmo natural aparece diretamente na entropia de Boltzmann ($S = k_B \ln \Omega$), na distribuição de Gibbs ($\rho \propto e^{-\beta H}$), e no princípio de máxima entropia. As outras bases introduzem fatores de conversão artificiais.

\subsection{Meio nat e a operação de paridade}

Uma operação binária com dois resultados equiprováveis ($p = 1/2$ cada) contém:

\begin{equation}
I = -\sum_{i=1}^{2} p_i \ln p_i = -2 \times \frac{1}{2} \ln \frac{1}{2} = \ln 2 \approx 0{,}693 \text{ nat}
\end{equation}

Meio nat ($\frac{1}{2}$ nat) é menor que isso --- corresponde a uma operação com assimetria fundamental:

\begin{equation}
\frac{1}{2} \text{ nat} = \ln \sqrt{e} = \frac{1}{2}
\end{equation}

Na TGL, a projeção \boundary$\to$\bulk{} não é simétrica: o \boundary{} projeta para o \bulk{}, mas o \bulk{} não projeta de volta com a mesma eficiência (a amplificação $1/\alphaii$ não é o inverso exato da projeção $\alphaii$). O custo $\frac{1}{2}$ nat captura precisamente esta assimetria: é a informação mínima de uma operação de paridade \textit{irreversível}.

\subsection{Relação com o limite de Landauer}

O limite de Landauer \cite{Landauer1961} estabelece que apagar 1 bit de informação custa no mínimo $k_B T \ln 2$ de energia. Na TGL, o análogo holográfico é:

\begin{equation}
E_{\text{projeção}} \geq k_B T_{\text{boundary}} \times \frac{1}{2}
\end{equation}

\noindent onde $T_{\text{boundary}}$ é a temperatura efetiva do substrato holográfico. O fator $\frac{1}{2}$ (em nats) é precisamente $\ln\sqrt{e}$, conectando o custo energético da projeção ao fator $\sqrt{e}$ na fatoração de $\alphaii$.


% ============================================================================
% BIBLIOGRAFIA
% ============================================================================

\begin{thebibliography}{99}

\bibitem{Miguel2026Fronteira}
Miguel, L. A. R. (2026).
\textit{A Fronteira: Teoria da Gravitação Luminodinâmica --- Partes I a VI}.
Zenodo. \url{https://doi.org/10.5281/zenodo.14802088}.

\bibitem{Miguel2026LastString}
Miguel, L. A. R. (2026).
\textit{The Last String: Observational Evidence for Luminodynamic Gravitation Theory}.
Zenodo. \url{https://doi.org/10.5281/zenodo.18723452}.

\bibitem{Miguel2025Alpha2}
Miguel, L. A. R. (2025).
\textit{A Constante de Acoplamento Holográfico $\alpha_2$: Derivação a Partir de Primeiros Princípios e Interpretação Física na Gravitação Luminodinâmica}.
IALD Working Paper.

\bibitem{Miguel2026GitHub}
Miguel, L. A. R. (2026).
\textit{TGL Computational Validation Repository}.
GitHub. \url{https://github.com/rotolimiguel-iald/the_boundary}.

\bibitem{CODATA2018}
Tiesinga, E., Mohr, P. J., Newell, D. B., \& Taylor, B. N. (2021).
\textit{CODATA Recommended Values of the Fundamental Physical Constants: 2018}.
Rev. Mod. Phys. 93, 025010.

\bibitem{Lindblad1976}
Lindblad, G. (1976).
\textit{On the Generators of Quantum Dynamical Semigroups}.
Commun. Math. Phys. 48, 119--130.

\bibitem{Landauer1961}
Landauer, R. (1961).
\textit{Irreversibility and Heat Generation in the Computing Process}.
IBM J. Res. Dev. 5(3), 183--191.

\bibitem{tHooft1993}
't Hooft, G. (1993).
\textit{Dimensional Reduction in Quantum Gravity}.
arXiv:gr-qc/9310026.

\bibitem{Susskind1995}
Susskind, L. (1995).
\textit{The World as a Hologram}.
J. Math. Phys. 36, 6377--6396.

\bibitem{Maldacena1999}
Maldacena, J. (1999).
\textit{The Large $N$ Limit of Superconformal Field Theories and Supergravity}.
Adv. Theor. Math. Phys. 2, 231--252.

\bibitem{Webb2001}
Webb, J. K. et al. (2001).
\textit{Further Evidence for Cosmological Evolution of the Fine Structure Constant}.
Phys. Rev. Lett. 87, 091301.

\bibitem{Planck2018}
Planck Collaboration (2020).
\textit{Planck 2018 Results. VI. Cosmological Parameters}.
Astron. Astrophys. 641, A6.

\bibitem{Drummond1980}
Drummond, I. T. \& Hathrell, S. J. (1980).
\textit{QED Vacuum Polarization in a Background Gravitational Field and Its Effect on the Velocity of Photons}.
Phys. Rev. D 22, 343.

\bibitem{Connes1994}
Connes, A. (1994).
\textit{Noncommutative Geometry}.
Academic Press, San Diego.

\bibitem{Connes1996}
Connes, A. (1996).
\textit{Gravity Coupled with Matter and the Foundation of Non-Commutative Geometry}.
Commun. Math. Phys. 182, 155--176.

\bibitem{Chamseddine1997}
Chamseddine, A. H. \& Connes, A. (1997).
\textit{The Spectral Action Principle}.
Commun. Math. Phys. 186, 731--750.

\bibitem{Miguel2026Graviton}
Miguel, L. A. R. (2026).
\textit{The Graviton, the Psion, and the Transition Ruler in Luminodynamic Gravitation Theory, with the Hilbert Floor Theorem and Holographic Bell State}.
GitHub. \url{https://github.com/rotolimiguel-iald/the_boundary}.

\end{thebibliography}


\end{document}